% QMB_n04_bell_theorem.tex — Kapitel 4: Bells Theorem
% Quellen: Dok. 023, 023a, 074
\chapter{Bells Theorem}
\label{ch:bell}

\section{Die Grundfrage}

Bell-Tests prüfen, ob die Korrelationen verschränkter Teilchen durch
eine lokal-realistische Theorie erklärt werden können — eine Theorie,
in der jedes Teilchen von Anfang an definierte Eigenschaften besitzt,
die unabhängig vom anderen Teilchen gemessen werden.

John Bell zeigte 1964, dass keine solche Theorie die
Quantenvorhersagen reproduzieren kann.\status{E} Das Argument ist
mathematisch scharf und lässt sich nicht durch technische Verfeinerungen
umgehen.

\section{Die CHSH-Ungleichung}

In der von Clauser, Horne, Shimony und Holt formulierten Form
(CHSH, 1969) lautet die Bell-Ungleichung:\status{E}
\begin{equation}
  |E(a,b)-E(a,c)| + |E(a',b)+E(a',c)| \le 2,
  \label{eq:chsh}
\end{equation}
wobei $E(a,b)$ die Korrelation der Messergebnisse bei den
Detektorwinkeln $a$ und $b$ bezeichnet.

Die Quantenmechanik verletzt diese Schranke bis zur
Tsirelson-Schranke $2\sqrt2\approx2{,}828$.\status{E} Alle
schlupflochfreien Experimente seit 2015 bestätigen die
QM-Vorhersage und widerlegen lokal-realistische Theorien.\status{E}

\section{Die drei Annahmen}

Bells Beweis setzt drei Annahmen voraus:
\begin{enumerate}
  \item \textbf{Lokalität:} Das Ergebnis bei Detektor A hängt nicht
        von der Einstellung bei Detektor B ab.
  \item \textbf{Realismus:} Die Messergebnisse sind vorher bestimmt
        (durch verborgene Variablen).
  \item \textbf{Messfreiheit:} Die Detektoreinstellungen werden
        unabhängig von den verborgenen Variablen gewählt.
\end{enumerate}
Mindestens eine dieser Annahmen ist falsch — das ist Bells
Leistung. Die QM gibt Lokalität und Realismus gemeinsam auf.

\section{Historische FFGFT-Ansätze und ihr heutiger Stand}

Die frühen FFGFT-Dokumente (Dok.~023, 023a, 074, Herbst 2025)
versuchten verschiedene Wege, die Verletzung mit lokaler Realität
zu versöhnen:

\begin{itemize}
  \item \textbf{Zeitfeld-Dämpfung} (Dok.~023): Die Korrelationen
        werden durch einen $\xi$-abhängigen Faktor abgeschwächt.
        Ergebnis: CHSH-Reduktion von $2{,}828$ auf $2{,}827$ —
        zu klein für heutige Experimente.\status{S}
  \item \textbf{Superdeterminismus} (Dok.~074): Die Messfreiheit
        ist verletzt, weil das Energiefeld Detektoreinstellungen
        und Teilcheneigenschaften gemeinsam bestimmt. Dieser Weg
        ist durch Dok.~230 überholt (s. Kapitel~\ref{ch:bell_ffgft}).
  \item \textbf{Topologische Verbindung} (Dok.~230, Mai 2026):
        Kein Zugeständnis an die Annahmen nötig — der gemeinsame
        Träger macht die Korrelation zu einer geometrischen
        Eigenschaft. Dieser Weg ist heute maßgeblich.\status{B}
\end{itemize}

\quelle{Dok.~023, 023a, 074, 190 (Präz.~11, 12), 230}
